% Chapter file for a thesis: the main file loads it with % Chapter file for a thesis: the main file loads it with % Chapter file for a thesis: the main file loads it with % Chapter file for a thesis: the main file loads it with \include{thesis-chapter}.
% The main file's preamble needs amsmath, graphicx, booktabs, hyperref and cleveref.

\chapter{Methodology}
\label{ch:methodology}

This chapter describes the study design, the data and the analysis.
\Cref{sec:design} sets out the design, \cref{sec:data} the data and
\cref{sec:analysis} the model.

\section{Study design}
\label{sec:design}

Describe the design and justify it with the literature~\cite{surname2023keyword}.

\section{Data}
\label{sec:data}

\Cref{tab:sample} summarises the sample.

\begin{table}[htbp]
  \centering
  \caption{Sample by faculty (example values: replace them)}
  \label{tab:sample}
  \begin{tabular}{@{}lrr@{}}
    \toprule
    Faculty     & Respondents & Share (\%) \\
    \midrule
    Science     & 120         & 40.0 \\
    Engineering & 105         & 35.0 \\
    Arts        & 75          & 25.0 \\
    \midrule
    Total       & 300         & 100.0 \\
    \bottomrule
  \end{tabular}
\end{table}

\section{Analysis}
\label{sec:analysis}

The outcome is modelled with logistic regression:
\begin{equation}
  \log \frac{p_i}{1 - p_i} = \beta_0 + \sum_{j=1}^{k} \beta_j x_{ij},
  \label{eq:logit}
\end{equation}
where $p_i$ is the probability of the outcome for respondent~$i$.
The coefficients in~\eqref{eq:logit} are estimated by maximum likelihood~\cite{author2020book}.
\Cref{fig:framework} shows the conceptual framework.

\begin{figure}[htbp]
  \centering
  \includegraphics[width=0.8\linewidth]{example-image}% TODO: replace with the real figure file
  \caption{Conceptual framework of the study.}
  \label{fig:framework}
\end{figure}

\section{Summary}
\label{sec:summary}

Summarise the chapter in one paragraph and link it to the next chapter.
.
% The main file's preamble needs amsmath, graphicx, booktabs, hyperref and cleveref.

\chapter{Methodology}
\label{ch:methodology}

This chapter describes the study design, the data and the analysis.
\Cref{sec:design} sets out the design, \cref{sec:data} the data and
\cref{sec:analysis} the model.

\section{Study design}
\label{sec:design}

Describe the design and justify it with the literature~\cite{surname2023keyword}.

\section{Data}
\label{sec:data}

\Cref{tab:sample} summarises the sample.

\begin{table}[htbp]
  \centering
  \caption{Sample by faculty (example values: replace them)}
  \label{tab:sample}
  \begin{tabular}{@{}lrr@{}}
    \toprule
    Faculty     & Respondents & Share (\%) \\
    \midrule
    Science     & 120         & 40.0 \\
    Engineering & 105         & 35.0 \\
    Arts        & 75          & 25.0 \\
    \midrule
    Total       & 300         & 100.0 \\
    \bottomrule
  \end{tabular}
\end{table}

\section{Analysis}
\label{sec:analysis}

The outcome is modelled with logistic regression:
\begin{equation}
  \log \frac{p_i}{1 - p_i} = \beta_0 + \sum_{j=1}^{k} \beta_j x_{ij},
  \label{eq:logit}
\end{equation}
where $p_i$ is the probability of the outcome for respondent~$i$.
The coefficients in~\eqref{eq:logit} are estimated by maximum likelihood~\cite{author2020book}.
\Cref{fig:framework} shows the conceptual framework.

\begin{figure}[htbp]
  \centering
  \includegraphics[width=0.8\linewidth]{example-image}% TODO: replace with the real figure file
  \caption{Conceptual framework of the study.}
  \label{fig:framework}
\end{figure}

\section{Summary}
\label{sec:summary}

Summarise the chapter in one paragraph and link it to the next chapter.
.
% The main file's preamble needs amsmath, graphicx, booktabs, hyperref and cleveref.

\chapter{Methodology}
\label{ch:methodology}

This chapter describes the study design, the data and the analysis.
\Cref{sec:design} sets out the design, \cref{sec:data} the data and
\cref{sec:analysis} the model.

\section{Study design}
\label{sec:design}

Describe the design and justify it with the literature~\cite{surname2023keyword}.

\section{Data}
\label{sec:data}

\Cref{tab:sample} summarises the sample.

\begin{table}[htbp]
  \centering
  \caption{Sample by faculty (example values: replace them)}
  \label{tab:sample}
  \begin{tabular}{@{}lrr@{}}
    \toprule
    Faculty     & Respondents & Share (\%) \\
    \midrule
    Science     & 120         & 40.0 \\
    Engineering & 105         & 35.0 \\
    Arts        & 75          & 25.0 \\
    \midrule
    Total       & 300         & 100.0 \\
    \bottomrule
  \end{tabular}
\end{table}

\section{Analysis}
\label{sec:analysis}

The outcome is modelled with logistic regression:
\begin{equation}
  \log \frac{p_i}{1 - p_i} = \beta_0 + \sum_{j=1}^{k} \beta_j x_{ij},
  \label{eq:logit}
\end{equation}
where $p_i$ is the probability of the outcome for respondent~$i$.
The coefficients in~\eqref{eq:logit} are estimated by maximum likelihood~\cite{author2020book}.
\Cref{fig:framework} shows the conceptual framework.

\begin{figure}[htbp]
  \centering
  \includegraphics[width=0.8\linewidth]{example-image}% TODO: replace with the real figure file
  \caption{Conceptual framework of the study.}
  \label{fig:framework}
\end{figure}

\section{Summary}
\label{sec:summary}

Summarise the chapter in one paragraph and link it to the next chapter.
.
% The main file's preamble needs amsmath, graphicx, booktabs, hyperref and cleveref.

\chapter{Methodology}
\label{ch:methodology}

This chapter describes the study design, the data and the analysis.
\Cref{sec:design} sets out the design, \cref{sec:data} the data and
\cref{sec:analysis} the model.

\section{Study design}
\label{sec:design}

Describe the design and justify it with the literature~\cite{surname2023keyword}.

\section{Data}
\label{sec:data}

\Cref{tab:sample} summarises the sample.

\begin{table}[htbp]
  \centering
  \caption{Sample by faculty (example values: replace them)}
  \label{tab:sample}
  \begin{tabular}{@{}lrr@{}}
    \toprule
    Faculty     & Respondents & Share (\%) \\
    \midrule
    Science     & 120         & 40.0 \\
    Engineering & 105         & 35.0 \\
    Arts        & 75          & 25.0 \\
    \midrule
    Total       & 300         & 100.0 \\
    \bottomrule
  \end{tabular}
\end{table}

\section{Analysis}
\label{sec:analysis}

The outcome is modelled with logistic regression:
\begin{equation}
  \log \frac{p_i}{1 - p_i} = \beta_0 + \sum_{j=1}^{k} \beta_j x_{ij},
  \label{eq:logit}
\end{equation}
where $p_i$ is the probability of the outcome for respondent~$i$.
The coefficients in~\eqref{eq:logit} are estimated by maximum likelihood~\cite{author2020book}.
\Cref{fig:framework} shows the conceptual framework.

\begin{figure}[htbp]
  \centering
  \includegraphics[width=0.8\linewidth]{example-image}% TODO: replace with the real figure file
  \caption{Conceptual framework of the study.}
  \label{fig:framework}
\end{figure}

\section{Summary}
\label{sec:summary}

Summarise the chapter in one paragraph and link it to the next chapter.
